\documentclass{article}
\usepackage[utf8]{inputenc}
\usepackage{a4wide}

\begin{document}

\section*{หาเลขใกล้}

กำหนดอาร์เรย์ของเลขจำนวนเต็ม\texttt{nums}และเลขจำนวนเต็มอีกสองค่าคือ\texttt{indexDiff}และ\texttt{valueDiff}ให้หาคู่ดัชนี\texttt{(i,\ j)}โดยที่:



1.\texttt{i\ !=\ j}



2.\texttt{abs(i\ -\ j)\ \textless{}=\ indexDiff}



3.\texttt{abs(nums{[}i{]}\ -\ nums{[}j{]})\ \textless{}=\ valueDiff}



ให้ตอบ\texttt{Y}\ul{หากมี}คู่ดัชนีดังกล่าว และตอบ\texttt{N}หาก\ul{ไม่มี}\\



\ul{\hfill\break

}



{Example 1:}



\begin{quote}

\textbf{input:}



4



1 2 3 1



30



\textbf{output:}



Y



\textbf{Explanation:}



สามารถหา\texttt{(i,\ j)\ =\ (0,\ 3)}ที่ทำให้ตรงเงื่อนไขทั้งสามข้อได้พร้อมกัน



1.\texttt{i\ !=\ j\ -\/-\textgreater{}\ 0\ !=\ 3}



2.\texttt{abs(i\ -\ j)\ \textless{}=\ indexDiff\ -\/-\textgreater{}\ abs(0\ -\ 3)\ \textless{}=\ 3}



3.\texttt{abs(nums{[}i{]}\ -\ nums{[}j{]})\ \textless{}=\ valueDiff\ -\/-\textgreater{}\ abs(1\ -\ 1)\ \textless{}=\ 0}

\end{quote}



{Example 2:}



\begin{quote}

\textbf{input:}



6



1 5 9 1 5 9



2 3



\textbf{output:}



N



\textbf{Explanation:}



ไม่สามารถหาคู่ดัชนี\texttt{(i,\ j)}ที่ทำให้ตรงเงื่อนไขทั้งสามข้อได้พร้อมกัน

\end{quote}



\subsection*{Input}
{บรรทัดที่ 1 ขนาดของอาร์เรย์{\texttt{nums}}{0 \textless{} nums \textless{}

2,000}}\\



บรรทัดที่ 2 ชุดของเลขจำนวนเต็มที่เว้นด้วยช่องว่าง 1 ช่อง ~ ~0 \textless{}

numsi\textless{} 1,005,000



บรรทัดที่ 3

เป็นเลขจำนวนเต็มสำหรับ\texttt{indexDiff}และ\texttt{valueDiff}{เว้นด้วยช่องว่าง

1 ช่อง ~0 \textless{} indexDiff, ValueDiff \textless{} 50}



\subsection*{Output}
อักขระ\texttt{Y}หรือ\texttt{N}\\



\end{document}
